% Options for packages loaded elsewhere
\PassOptionsToPackage{unicode}{hyperref}
\PassOptionsToPackage{hyphens}{url}
\documentclass[
]{article}
\usepackage{xcolor}
\usepackage[a4paper,margin=2cm]{geometry}
\usepackage{amsmath,amssymb}
\setcounter{secnumdepth}{-\maxdimen} % remove section numbering
\usepackage{iftex}
\ifPDFTeX
  \usepackage[T1]{fontenc}
  \usepackage[utf8]{inputenc}
  \usepackage{textcomp} % provide euro and other symbols
\else % if luatex or xetex
  \usepackage{unicode-math} % this also loads fontspec
  \defaultfontfeatures{Scale=MatchLowercase}
  \defaultfontfeatures[\rmfamily]{Ligatures=TeX,Scale=1}
\fi
\usepackage{lmodern}
\ifPDFTeX\else
  % xetex/luatex font selection
\fi
% Use upquote if available, for straight quotes in verbatim environments
\IfFileExists{upquote.sty}{\usepackage{upquote}}{}
\IfFileExists{microtype.sty}{% use microtype if available
  \usepackage[]{microtype}
  \UseMicrotypeSet[protrusion]{basicmath} % disable protrusion for tt fonts
}{}
\makeatletter
\@ifundefined{KOMAClassName}{% if non-KOMA class
  \IfFileExists{parskip.sty}{%
    \usepackage{parskip}
  }{% else
    \setlength{\parindent}{0pt}
    \setlength{\parskip}{6pt plus 2pt minus 1pt}}
}{% if KOMA class
  \KOMAoptions{parskip=half}}
\makeatother
\usepackage{color}
\usepackage{fancyvrb}
\newcommand{\VerbBar}{|}
\newcommand{\VERB}{\Verb[commandchars=\\\{\}]}
\DefineVerbatimEnvironment{Highlighting}{Verbatim}{commandchars=\\\{\}}
% Add ',fontsize=\small' for more characters per line
\newenvironment{Shaded}{}{}
\newcommand{\AlertTok}[1]{\textcolor[rgb]{1.00,0.00,0.00}{\textbf{#1}}}
\newcommand{\AnnotationTok}[1]{\textcolor[rgb]{0.38,0.63,0.69}{\textbf{\textit{#1}}}}
\newcommand{\AttributeTok}[1]{\textcolor[rgb]{0.49,0.56,0.16}{#1}}
\newcommand{\BaseNTok}[1]{\textcolor[rgb]{0.25,0.63,0.44}{#1}}
\newcommand{\BuiltInTok}[1]{\textcolor[rgb]{0.00,0.50,0.00}{#1}}
\newcommand{\CharTok}[1]{\textcolor[rgb]{0.25,0.44,0.63}{#1}}
\newcommand{\CommentTok}[1]{\textcolor[rgb]{0.38,0.63,0.69}{\textit{#1}}}
\newcommand{\CommentVarTok}[1]{\textcolor[rgb]{0.38,0.63,0.69}{\textbf{\textit{#1}}}}
\newcommand{\ConstantTok}[1]{\textcolor[rgb]{0.53,0.00,0.00}{#1}}
\newcommand{\ControlFlowTok}[1]{\textcolor[rgb]{0.00,0.44,0.13}{\textbf{#1}}}
\newcommand{\DataTypeTok}[1]{\textcolor[rgb]{0.56,0.13,0.00}{#1}}
\newcommand{\DecValTok}[1]{\textcolor[rgb]{0.25,0.63,0.44}{#1}}
\newcommand{\DocumentationTok}[1]{\textcolor[rgb]{0.73,0.13,0.13}{\textit{#1}}}
\newcommand{\ErrorTok}[1]{\textcolor[rgb]{1.00,0.00,0.00}{\textbf{#1}}}
\newcommand{\ExtensionTok}[1]{#1}
\newcommand{\FloatTok}[1]{\textcolor[rgb]{0.25,0.63,0.44}{#1}}
\newcommand{\FunctionTok}[1]{\textcolor[rgb]{0.02,0.16,0.49}{#1}}
\newcommand{\ImportTok}[1]{\textcolor[rgb]{0.00,0.50,0.00}{\textbf{#1}}}
\newcommand{\InformationTok}[1]{\textcolor[rgb]{0.38,0.63,0.69}{\textbf{\textit{#1}}}}
\newcommand{\KeywordTok}[1]{\textcolor[rgb]{0.00,0.44,0.13}{\textbf{#1}}}
\newcommand{\NormalTok}[1]{#1}
\newcommand{\OperatorTok}[1]{\textcolor[rgb]{0.40,0.40,0.40}{#1}}
\newcommand{\OtherTok}[1]{\textcolor[rgb]{0.00,0.44,0.13}{#1}}
\newcommand{\PreprocessorTok}[1]{\textcolor[rgb]{0.74,0.48,0.00}{#1}}
\newcommand{\RegionMarkerTok}[1]{#1}
\newcommand{\SpecialCharTok}[1]{\textcolor[rgb]{0.25,0.44,0.63}{#1}}
\newcommand{\SpecialStringTok}[1]{\textcolor[rgb]{0.73,0.40,0.53}{#1}}
\newcommand{\StringTok}[1]{\textcolor[rgb]{0.25,0.44,0.63}{#1}}
\newcommand{\VariableTok}[1]{\textcolor[rgb]{0.10,0.09,0.49}{#1}}
\newcommand{\VerbatimStringTok}[1]{\textcolor[rgb]{0.25,0.44,0.63}{#1}}
\newcommand{\WarningTok}[1]{\textcolor[rgb]{0.38,0.63,0.69}{\textbf{\textit{#1}}}}
\usepackage{longtable,booktabs,array}
\usepackage{caption}
\captionsetup[table]{skip=6pt}
\newcounter{none} % for unnumbered tables
\usepackage{calc} % for calculating minipage widths
% Correct order of tables after \paragraph or \subparagraph
\usepackage{etoolbox}
\makeatletter
\patchcmd\longtable{\par}{\if@noskipsec\mbox{}\fi\par}{}{}
\makeatother
% Allow footnotes in longtable head/foot
\IfFileExists{footnotehyper.sty}{\usepackage{footnotehyper}}{\usepackage{footnote}}
\makesavenoteenv{longtable}
\setlength{\emergencystretch}{3em} % prevent overfull lines
\providecommand{\tightlist}{%
  \setlength{\itemsep}{0pt}\setlength{\parskip}{0pt}}
\usepackage{bookmark}
\IfFileExists{xurl.sty}{\usepackage{xurl}}{} % add URL line breaks if available
\urlstyle{same}
% fallback for those not using the hyperref driver hyperxmp:
\makeatletter
\@ifundefined{xmpquote}{\newcommand{\xmpquote}[1]{#1}}{}
\makeatother
\hypersetup{
  hidelinks,
  pdfcreator={LaTeX via pandoc}}

\author{}
\date{}

\begin{document}

\textbf{Principal Component Analysis vs Factor Analysis}

MA2003B - Application of Multivariate Methods in Data Science

MA2003B Course Team

Tecnológico de Monterrey

13:59:05 2026-08-28

\begin{center}\rule{0.5\linewidth}{0.5pt}\end{center}

\section{Introduction to Principal Component Analysis and Factor
Analysis}

\subsection{Why Dimensionality Reduction?}

Dimensionality reduction is a fundamental technique in multivariate
analysis that seeks to:

\begin{itemize}
\item
  Handle high-dimensional data effectively
\item
  Reduce computational complexity
\item
  Remove noise and redundancy
\item
  Improve model interpretability
\item
  Enable data visualization
\end{itemize}

\subsection{Two Main Approaches}

\begin{itemize}
\item
  \textbf{Principal Component Analysis (PCA)}: Data-driven approach,
  variance maximization
\item
  \textbf{Factor Analysis (FA)}: Model-based approach, identification of
  latent constructs
\end{itemize}

\section{Principal Component Analysis (PCA)}

\subsection{What is PCA?}

PCA is a dimensionality reduction technique that transforms correlated
variables into uncorrelated principal components, maximizing the
variance along each new axis and ordering the components by explained
variance.

\subsection{Mathematical Foundation}

Given a data matrix \(\mathbf{X}\) with \(n\) observations and \(p\)
variables:

\begin{enumerate}
\item
  Center the data:
  \(\mathbf{X}_{\text{centered }} = \mathbf{X} - \overline{\mathbf{X}}\)
\item
  Calculate the covariance matrix:
  \(\mathbf{S} = \left( \frac{1}{n - 1} \right)\mathbf{X}_{\text{centered}}^{T}\mathbf{X}_{\text{centered}}\)
\item
  Find eigenvalues \(\lambda_{i}\) and eigenvectors \(\mathbf{v}_{i}\)
  of \(\mathbf{S}\)
\item
  Principal components:
  \(\mathbf{\text{PC}}_{i} = \mathbf{X}_{\text{centered }}\mathbf{v}_{i}\)
\end{enumerate}

\subsection{PCA Example}

\begin{Shaded}
\begin{Highlighting}[]
\ImportTok{import}\NormalTok{ numpy }\ImportTok{as}\NormalTok{ np}
\ImportTok{from}\NormalTok{ sklearn.decomposition }\ImportTok{import}\NormalTok{ PCA}

\CommentTok{\# Simple 3x2 data matrix}
\NormalTok{X }\OperatorTok{=}\NormalTok{ np.array([[}\DecValTok{5}\NormalTok{, }\DecValTok{3}\NormalTok{],}
\NormalTok{              [}\DecValTok{3}\NormalTok{, }\DecValTok{1}\NormalTok{],}
\NormalTok{              [}\DecValTok{1}\NormalTok{, }\DecValTok{3}\NormalTok{]])}

\CommentTok{\# Apply PCA}
\NormalTok{pca }\OperatorTok{=}\NormalTok{ PCA()}
\NormalTok{X\_transformed }\OperatorTok{=}\NormalTok{ pca.fit\_transform(X)}

\CommentTok{\# Results}
\NormalTok{eigenvalues }\OperatorTok{=}\NormalTok{ pca.explained\_variance\_}
\NormalTok{variance\_ratio }\OperatorTok{=}\NormalTok{ pca.explained\_variance\_ratio\_}

\BuiltInTok{print}\NormalTok{(}\SpecialStringTok{f"Eigenvalues: }\SpecialCharTok{\{}\NormalTok{eigenvalues}\SpecialCharTok{\}}\SpecialStringTok{"}\NormalTok{)}
\BuiltInTok{print}\NormalTok{(}\SpecialStringTok{f"PC1 explains }\SpecialCharTok{\{}\NormalTok{variance\_ratio[}\DecValTok{0}\NormalTok{]}\SpecialCharTok{:.1\%\}}\SpecialStringTok{ of variance"}\NormalTok{)}
\end{Highlighting}
\end{Shaded}

\subsection{Key Results}

\begin{itemize}
\item
  PC1 explains the most variance (highest eigenvalue)
\item
  Components are uncorrelated by construction
\item
  Original data can be reconstructed from the components
\end{itemize}

\subsection{Component Retention Criteria}

\begin{itemize}
\item
  \textbf{Kaiser Criterion}: Keep components with \(\lambda_{i} > 1\)
\item
  \textbf{Scree Plot}: Look for the ``elbow'' in the eigenvalue plot
\item
  \textbf{Cumulative Variance}: Retain enough components for the desired
  variance (e.g., 80\%)
\item
  \textbf{Parallel Analysis}: Compare with eigenvalues of random data
\end{itemize}

\section{Factor Analysis}

\subsection{What is Factor Analysis?}

Factor analysis assumes that observed variables are linear combinations
of common factors (shared latent constructs) and unique factors
(variable-specific variance).

\subsection{The Common Factor Model}

For each observed variable \(x_{j}\):

\(x_{j} = \mu_{j} + \sum_{i = 1}^{m}\lambda_{ji}f_{i} + \varepsilon_{j}\)

Where:

\begin{itemize}
\item
  \(\lambda_{ji}\): Factor loading (correlation between \(x_{j}\) and
  \(f_{i}\))
\item
  \(f_{i}\): Common factor (latent variable)
\item
  \(\varepsilon_{j}\): Unique factor (error term)
\end{itemize}

\subsection{Factor Analysis Example}

\begin{Shaded}
\begin{Highlighting}[]
\ImportTok{import}\NormalTok{ numpy }\ImportTok{as}\NormalTok{ np}
\ImportTok{from}\NormalTok{ factor\_analyzer }\ImportTok{import}\NormalTok{ FactorAnalyzer}

\CommentTok{\# Correlation matrix}
\NormalTok{R }\OperatorTok{=}\NormalTok{ np.array([[}\FloatTok{1.00}\NormalTok{, }\FloatTok{0.60}\NormalTok{, }\FloatTok{0.48}\NormalTok{],}
\NormalTok{              [}\FloatTok{0.60}\NormalTok{, }\FloatTok{1.00}\NormalTok{, }\FloatTok{0.72}\NormalTok{],}
\NormalTok{              [}\FloatTok{0.48}\NormalTok{, }\FloatTok{0.72}\NormalTok{, }\FloatTok{1.00}\NormalTok{]])}

\CommentTok{\# Perform Factor Analysis}
\NormalTok{fa }\OperatorTok{=}\NormalTok{ FactorAnalyzer(n\_factors}\OperatorTok{=}\DecValTok{1}\NormalTok{, rotation}\OperatorTok{=}\VariableTok{None}\NormalTok{)}
\NormalTok{fa.fit(R)}

\CommentTok{\# Results}
\NormalTok{loadings }\OperatorTok{=}\NormalTok{ fa.loadings\_}
\NormalTok{communalities }\OperatorTok{=}\NormalTok{ fa.get\_communalities()}
\NormalTok{uniqueness }\OperatorTok{=}\NormalTok{ fa.get\_uniquenesses()}

\BuiltInTok{print}\NormalTok{(}\SpecialStringTok{f"Factor loadings:}\CharTok{\textbackslash{}\textbackslash{}}\SpecialStringTok{n}\SpecialCharTok{\{}\NormalTok{loadings}\SpecialCharTok{\}}\SpecialStringTok{"}\NormalTok{)}
\BuiltInTok{print}\NormalTok{(}\SpecialStringTok{f"Communalities: }\SpecialCharTok{\{}\NormalTok{communalities}\SpecialCharTok{\}}\SpecialStringTok{"}\NormalTok{)}
\BuiltInTok{print}\NormalTok{(}\SpecialStringTok{f"Uniquenesses: }\SpecialCharTok{\{}\NormalTok{uniqueness}\SpecialCharTok{\}}\SpecialStringTok{"}\NormalTok{)}
\end{Highlighting}
\end{Shaded}

\subsection{Key Concepts}

\begin{itemize}
\item
  \textbf{Loadings}: Correlations between variables and factors
\item
  \textbf{Communalities}: Variance explained by common factors
\item
  \textbf{Uniquenesses}: Variable-specific variance
\end{itemize}

\subsection{Factor Rotation}

Factor rotation is a technique applied after factor extraction to make
factors easier to interpret. It ``rotates'' the factor axes to achieve
\textbf{simple structure}: each variable should load highly on only one
factor and low on others.

\textbf{The problem:} Initial extraction often produces factors where
many variables have moderate loadings on multiple factors, making
interpretation difficult.

\textbf{The solution:} Rotation redistributes variance among factors
without changing the total variance explained.

\textbf{Two main types:}

\begin{itemize}
\item
  \textbf{Orthogonal rotation} (e.g., Varimax): Keeps factors
  uncorrelated (independent)

  \begin{itemize}
  \item
    Simpler interpretation, but may be unrealistic if constructs
    actually correlate
  \end{itemize}
\item
  \textbf{Oblique rotation} (e.g., Promax, Oblimin): Allows factors to
  correlate

  \begin{itemize}
  \item
    More realistic for real-world data, slightly more complex
    interpretation
  \end{itemize}
\end{itemize}

{\def\LTcaptype{none} % do not increment counter
\begin{longtable}[]{@{}cccc@{}}
\toprule\noalign{}
\textbf{Method} & \textbf{Description} & \textbf{Assumption} &
\textbf{Use Case} \\
\midrule\noalign{}
\endhead
\bottomrule\noalign{}
\endlastfoot
Varimax & Orthogonal rotation & Uncorrelated factors & Simple
structure \\
Promax & Oblique rotation & Factors may correlate & Realistic models \\
Quartimax & Simplify variables & Balance factors & General use \\
Oblimin & Oblique rotation & Flexible correlation & Complex data \\
\end{longtable}
}

\subsection{Why Rotate?}

\begin{itemize}
\item
  Improve the interpretability of factor loadings
\item
  Achieve ``simple structure'' (variables that load highly on few
  factors)
\item
  Different rotations can reveal different substantive interpretations
\end{itemize}

\section{Comparison: PCA vs Factor Analysis}

\subsection{Key Differences}

{\def\LTcaptype{none} % do not increment counter
\begin{longtable}[]{@{}ccc@{}}
\toprule\noalign{}
\textbf{Aspect} & \textbf{PCA} & \textbf{Factor Analysis} \\
\midrule\noalign{}
\endhead
\bottomrule\noalign{}
\endlastfoot
Objective & Variance maximization & Latent construct identification \\
Model & Data-based & Theory-based \\
Components/Factors & All variance & Only common variance \\
Rotation & Not typically used & Essential for interpretation \\
Assumptions & Minimal & Multivariate normality \\
Estimation & Eigenvalue decomposition & Maximum likelihood/PAF \\
Output & Principal components & Factor loadings \\
\end{longtable}
}

\subsection{When to Use Each Method}

\subsubsection{Use PCA when:}

\begin{itemize}
\item
  Data reduction is the main goal
\item
  No theoretical model exists
\item
  All variance is of interest
\item
  Prediction is the objective
\item
  Data visualization is needed
\end{itemize}

\subsubsection{Use Factor Analysis when:}

\begin{itemize}
\item
  Identifying latent constructs
\item
  The analysis is theory-driven
\item
  Developing a measurement model
\item
  Understanding relationships between variables
\item
  Performing scale validation/development
\end{itemize}

\section{Complete Analysis Workflow}

Before performing Factor Analysis, we must verify that our data is
suitable:

\textbf{Kaiser-Meyer-Olkin (KMO) Test}

\begin{itemize}
\item
  Measures if your variables have enough in common to group into factors
\item
  High KMO means variables cluster well together; low KMO means they are
  too independent
\item
  Ranges from 0 to 1
\item
  Interpretation: \textgreater{} 0.9 excellent, \textgreater{} 0.8 good,
  \textgreater{} 0.6 acceptable, \textless{} 0.5 unacceptable
\item
  \textbf{What it tests:} Are correlations strong enough for FA to be
  useful?
\end{itemize}

\textbf{Bartlett's Test of Sphericity}

\begin{itemize}
\item
  Tests if the correlation matrix is significantly different from an
  identity matrix
\item
  Uses chi-square statistic
\item
  Interpretation: \(p < 0.05\) means variables are correlated (good for
  FA)
\item
  \textbf{What it tests:} Are variables correlated at all?
\end{itemize}

\textbf{Key difference:} Bartlett tests if correlations exist, KMO tests
if they are strong enough.

\begin{Shaded}
\begin{Highlighting}[]
\ImportTok{import}\NormalTok{ numpy }\ImportTok{as}\NormalTok{ np}
\ImportTok{from}\NormalTok{ sklearn.decomposition }\ImportTok{import}\NormalTok{ PCA}
\ImportTok{from}\NormalTok{ sklearn.preprocessing }\ImportTok{import}\NormalTok{ StandardScaler}
\ImportTok{from}\NormalTok{ factor\_analyzer }\ImportTok{import}\NormalTok{ FactorAnalyzer}
\ImportTok{from}\NormalTok{ factor\_analyzer.factor\_analyzer }\ImportTok{import}\NormalTok{ calculate\_kmo, calculate\_bartlett\_sphericity}

\CommentTok{\# 1. Load and prepare data}
\NormalTok{X }\OperatorTok{=}\NormalTok{ np.random.randn(}\DecValTok{100}\NormalTok{, }\DecValTok{5}\NormalTok{)  }\CommentTok{\# Your data here}

\CommentTok{\# 2. Standardize}
\NormalTok{scaler }\OperatorTok{=}\NormalTok{ StandardScaler()}
\NormalTok{X\_scaled }\OperatorTok{=}\NormalTok{ scaler.fit\_transform(X)}

\CommentTok{\# 3. Check suitability for FA}
\NormalTok{kmo\_all, kmo\_model }\OperatorTok{=}\NormalTok{ calculate\_kmo(X\_scaled)}
\NormalTok{chi\_square\_value, p\_value }\OperatorTok{=}\NormalTok{ calculate\_bartlett\_sphericity(X\_scaled)}

\BuiltInTok{print}\NormalTok{(}\SpecialStringTok{f"KMO: }\SpecialCharTok{\{}\NormalTok{kmo\_model}\SpecialCharTok{:.3f\}}\SpecialStringTok{ (\textgreater{}0.6 is good)"}\NormalTok{)}
\BuiltInTok{print}\NormalTok{(}\SpecialStringTok{f"Bartlett\textquotesingle{}s test p{-}value: }\SpecialCharTok{\{}\NormalTok{p\_value}\SpecialCharTok{:.3f\}}\SpecialStringTok{ (\textless{}0.05 is good)"}\NormalTok{)}

\CommentTok{\# 4. Determine number of factors}
\NormalTok{pca }\OperatorTok{=}\NormalTok{ PCA()}
\NormalTok{pca.fit(X\_scaled)}
\NormalTok{eigenvalues }\OperatorTok{=}\NormalTok{ pca.explained\_variance\_}
\NormalTok{n\_factors }\OperatorTok{=} \BuiltInTok{sum}\NormalTok{(eigenvalues }\OperatorTok{\textgreater{}} \DecValTok{1}\NormalTok{)  }\CommentTok{\# Kaiser criterion}

\CommentTok{\# 5. Perform Factor Analysis}
\NormalTok{fa }\OperatorTok{=}\NormalTok{ FactorAnalyzer(n\_factors}\OperatorTok{=}\NormalTok{n\_factors, rotation}\OperatorTok{=}\StringTok{\textquotesingle{}varimax\textquotesingle{}}\NormalTok{)}
\NormalTok{fa.fit(X\_scaled)}

\CommentTok{\# 6. Compare results}
\NormalTok{loadings }\OperatorTok{=}\NormalTok{ fa.loadings\_}
\NormalTok{communalities }\OperatorTok{=}\NormalTok{ fa.get\_communalities()}
\NormalTok{variance\_explained }\OperatorTok{=}\NormalTok{ fa.get\_factor\_variance()}
\end{Highlighting}
\end{Shaded}

\section{Applications and Best Practices}

\subsection{Real-World Applications}

\subsubsection{Finance}

\begin{itemize}
\item
  Risk factors in stock markets
\item
  Portfolio optimization
\item
  Credit scoring models
\end{itemize}

\subsubsection{Health}

\begin{itemize}
\item
  Patient satisfaction surveys
\item
  Symptom clustering
\item
  Quality of life measures
\end{itemize}

\subsubsection{Marketing}

\begin{itemize}
\item
  Customer segmentation
\item
  Brand perception studies
\item
  Product attribute analysis
\end{itemize}

\subsubsection{Social Sciences}

\begin{itemize}
\item
  Personality assessment
\item
  Attitude measurement
\item
  Educational testing
\end{itemize}

\subsection{Best Practices}

\subsubsection{Data Preparation}

\begin{itemize}
\item
  Ensure adequate sample size (5-10 observations per variable)
\item
  Check for multivariate normality
\item
  Handle missing data appropriately
\item
  Consider variable standardization
\end{itemize}

\subsubsection{Model Selection}

\begin{itemize}
\item
  Use KMO and Bartlett's tests for FA suitability
\item
  Compare multiple factor retention criteria
\item
  Consider both orthogonal and oblique rotations
\item
  Validate results with cross-validation
\end{itemize}

\subsubsection{Interpretation}

\begin{itemize}
\item
  Focus on the substantive meaning of the factors
\item
  Use factor loadings \textgreater{} 0.3 for interpretation
\item
  Consider correlations between factors in oblique rotations
\item
  Validate with external criteria when possible
\end{itemize}

\section{Conclusion}

\subsection{Key Takeaways}

\begin{itemize}
\item
  \textbf{PCA} and \textbf{Factor Analysis} serve different but
  complementary purposes
\item
  Choose the method based on research objectives and data
  characteristics
\item
  Always validate assumptions and interpret results substantively
\item
  Modern software makes implementation straightforward
\end{itemize}

\subsection{Next Steps}

\begin{itemize}
\item
  Practice with real datasets
\item
  Compare PCA and FA on the same data
\item
  Explore advanced techniques (confirmatory FA, structural equation
  modeling)
\item
  Apply to your own research questions
\end{itemize}

\section{References}

\subsection{Key References}

\begin{itemize}
\item
  \textbf{Fabrigar, L. R., \& Wegener, D. T.} (2011). Exploratory Factor
  Analysis. Oxford University Press.
\item
  \textbf{Hair, J. F., et al.} (2019). Multivariate Data Analysis.
  Cengage Learning.
\item
  \textbf{Jolliffe, I. T.} (2002). Principal Component Analysis.
  Springer.
\item
  \textbf{Tabachnick, B. G., \& Fidell, L. S.} (2013). Using
  Multivariate Statistics. Pearson.
\end{itemize}

\subsection{Software Resources}

\begin{itemize}
\item
  Python: \texttt{scikit-learn}, \texttt{factor-analyzer},
  \texttt{statsmodels}
\item
  R: \texttt{psych}, \texttt{FactoMineR}, \texttt{lavaan}
\item
  SPSS: Factor Analysis Module
\item
  SAS: PROC FACTOR
\end{itemize}

\subsection{Online Resources}

\begin{itemize}
\item
  UCLA Statistical Consulting Group
\item
  StatQuest YouTube channel
\item
  Towards Data Science articles
\end{itemize}

\end{document}

